% !PRINTED 447
\weberpage{447}
%<<CONT
achtz\"ahlig. Da er durch die Substitutionen 2\textsuperscript{ten} Grades $\omega$, $\Theta\omega$, $\Theta^2\omega$, $\Theta^3\omega$ unge\"andert bleibt, und da die Pole dieser Substitutionen auf $A$ liegen (weil $A$ durch sie unge\"andert bleibt), w\"ahrend $a$ nicht auf $A$ liegt, so m\"ussen die Hauptaxen dieser vier Substitutionen $\omega$, $\Theta\omega$, $\Theta^2\omega$, $\Theta^3\omega$ durch den Punkt $a$ gehen.

Die Pole der vier Substitutionen $\omega$, $\Theta\omega$, $\Theta^2\omega$, $\Theta^3\omega$, die wir $a_1$, $a_2$, $a_3$, $a_4$ nennen wollen, sind als Bilder des Punktes $a$ gleichfalls achtz\"ahlige Pole, und m\"ussen, wie $a$, Pole von Substitutionen vierter Ordnung sein, die jedenfalls alle von $\Theta$ und $\Theta^3$ verschieden sind, weil sonst $\Theta^2$ einer jener vier Substitutionen gleich sein m\"usste, was nicht der Fall ist. Durch $\omega$ bleiben nur zwei Punkte der Axe $A$, n\"amlich die Pole von $\omega$ und von $\Theta^2\omega$, unge\"andert.

Betrachten wir nun die beiden anderen Pole $c_1$, $c_2$ der Substitution $\Theta$. Diese Punkte liegen gleichfalls auf der Hauptaxe $A$, sind aber von den Punkten $a_1$, $a_2$, $a_3$, $a_4$ verschieden, weil diese, wie schon bemerkt, nicht unter den Polen von $\Theta$ vorkommen. Nun ist $\omega\Theta\omega = \Theta$, also

\[
\omega\,\Theta(x) = \Theta\,\omega(x),
\]
und wenn nun $(x) = \Theta(x)$ ist, so ist auch

\[
\omega(x) = \Theta\,\omega(x);
\]
d. h. wenn $(x)$ die Coordinaten eines Poles von $\Theta$ sind, so haben die $\omega(x)$ die gleiche Bedeutung. Durch die Transformation $\omega$ gehen also die Pole von $\Theta$ nur in einander \"uber. Der Pol $a$ bleibt unge\"andert durch $\omega$; $c_1$ und $c_2$ dagegen k\"onnen nicht unge\"andert bleiben, weil sie unter den Polen von $\omega$ nicht vorkommen, und m\"ussen also in einander \"ubergehen.

Wir haben also folgende Uebersicht:

\begin{center}
\begin{tabular}{@{}lllllll@{}}
Unge\"andert durch & $\omega$ und $\Theta^2\omega$ & auf & $A$ & nur die Punkte & $a_1, a_3$ & \\
\textquotedbl\ \textquotedbl & $\Theta\omega$\ \textquotedbl\ $\Theta^3\omega$ & \textquotedbl & $A$ & \textquotedbl\ \textquotedbl\ \textquotedbl & $a_2, a_4$ & \\
\textquotedbl\ \textquotedbl & $\Theta$\ \textquotedbl\ $\Theta^3$ & \textquotedbl & $A$ & \textquotedbl\ \textquotedbl\ \textquotedbl & $c_1, c_2.$ &
\end{tabular}
\end{center}

Ein von diesen sechs verschiedener Punkt $x$ von $A$ geht nur durch die Substitutionen 1 und $\Theta^2$ in sich selbst \"uber und kann daher ein \gesperrt{zweiz\"ahliger Pol} genannt werden. Dann gilt der Satz:

\noindent 3.\ \gesperrt{Ein zweiz\"ahliger Pol geht durch die aus $\omega$ und $\Theta$ abgeleitete Gruppe 8\textsuperscript{ten} Grades in vier verschiedene Lagen \"uber.}
